\bookchapter{第三十五章：兒童與手機——數字時代之教育指南}

\chapterintro{本章不講“手機乃洪水猛獸”，亦不鼓吹“越早接觸越聰明”。我等按年齡分段賜設備、設規則、教素養、護健康，一步步來。}

\booksection{一、幾歲賜孩子手機？}

無標準答案，但有年齡基線。過早賜手機，孩子大腦前額葉尚未發育成熟，自控力跟不上算法之“誘惑設計”；過晚賜，又或使他於同伴中“社交脫節”。下面乃分齡建議：

\begin{itemize}
  \item \textbf{3-6歲}：儘量不碰手機。如需使用，單次不超過15分鐘，內容以家長精選之教育動畫爲主，且須\textbf{陪同觀看}。此一階段屏幕對視力與留意力發育影響最大。
  \item \textbf{7-9歲}：可於家長設備上"借用"，用以網課、聽故事、查資料。開始建立"使用前申請、用完歸還"之規則，不配置專屬手機。
  \item \textbf{10-12歲}：可考慮賜一臺"功能機"或兒童手錶，解決上下學聯絡。如確需智能機，應開啓嚴格之家長控制。
  \item \textbf{13歲以上}：進入初中，社交需求增強，可賜智能機，但需配合"使用公約"與定期覆盤。
\end{itemize}

\begin{tcbinfo}{"夠用遂好"原則}
孩子之第一臺手機不須旗艦。處理器之“過剩性能”往往爲用來跑遊戲。選一臺電池夠大、性能中等、屏幕護眼之千元機或官方翻新機，反而更利於管控。預算參考：1500-2500元區間即可滿足絕大多數學生需求。

\end{tcbinfo}

\booksubsection{智能手錶作爲過渡方案}

若只欲解決“放學後能否聯繫上”之問題，兒童手錶比手機更合適。其只能打電話、發語音、定位，無瀏覽器與短視頻，天然限制矣“沉迷入口”：

\begin{itemize}
  \item \textbf{小天才電話手錶}：生態最完整，同齡人社交"碰一碰加好友"乃其核心賣點，但專屬App生態較封閉。
  \item \textbf{華爲兒童手錶}：與華爲手機聯動好，定位精準，支持視頻通話，性價比高。
  \item \textbf{小米兒童手錶}：價格親民，基礎功能齊全，適合預算有限之家庭。
\end{itemize}

買手錶時留意：選擇支持\textbf{4G全網通}之型號（2G/3G正逐步退網），並確認運營商兒童套餐資費。

\booksection{二、家長控制系統全攻略}

設備賜至孩子之前，第一件事非開機，而乃設置家長控制。兩大平臺皆有原生方案：

\booksubsection{iOS：屏幕使用時間}

\begin{itemize}
  \item \textbf{App限額}：設置 > 屏幕使用時間 > App限額，可爲每個App類別設定每日可用時長（如遊戲30分鐘）。
  \item \textbf{停用時間}：設定每日"熄燈時段"（如22:00-7:00），期間凡App變灰禁用，僅電話可用。
  \item \textbf{內容限制}：可按年齡段過濾影視、網頁、App安裝，並禁止App內購買。
  \item \textbf{家庭共享}：用家長之iPhone遠程管理孩子之設備，批准下載請求、查看使用報告。
\end{itemize}

\booksubsection{安卓：Google Family Link / 品牌家長控制}

\begin{itemize}
  \item \textbf{Google Family Link}：谷歌官方方案，可遠程管理設備使用時長、鎖定設備、查看活動報告。國內安卓需視機型可用性。
  \item \textbf{小米兒童空間}：設置 > 兒童空間，可設定使用時長、白名單App，退出需家長密碼。
  \item \textbf{華爲學生模式}：設置 > 健康使用手機 > 學生模式，限制使用時段與App，支持遠程管理。
  \item \textbf{OPPO/vivo兒童模式}：於"設置"或專屬"兒童空間"App中開啓，原理類似。
\end{itemize}

\begin{center}
\begin{tabularx}{\linewidth}{|X|X|X|}
\hline
\textbf{功能} & \textbf{iOS 屏幕使用時間} & \textbf{安卓 Family Link / 品牌方案} \tabularnewline
\hline
App限額 & 支持，按類別 & 支持，按App \tabularnewline
停用時間 & 支持 & 支持（品牌方案有差異） \tabularnewline
內容過濾 & 系統級，較徹底 & 依賴品牌，參差不齊 \tabularnewline
遠程管理 & 家庭共享 & Family Link / 品牌雲 \tabularnewline
安裝App審批 & 需家長批准 & Family Link 支持 \tabularnewline
\hline
\end{tabularx}
\end{center}

\begin{tcbwarn}{家長控制非"一勞永逸"}
孩子總能尋到繞過之方法（網上搜“何以破解屏幕使用時間”之教程一大把）。技術手段只爲底線，真正之防線乃親子溝通與信任。建議每月與孩子一起復盤使用報告，而非單方面監控。

\end{tcbwarn}

\booksection{三、青少年模式}

各大內容平臺皆內置矣“青少年模式”，開啓後將限制使用時長、過濾不良內容、關閉充值打賞。此乃國家法規要求之強制功能，但\textbf{須家長主動開啓並設置密碼}：

\begin{itemize}
  \item \textbf{抖音/快手青少年模式}：每天限40分鐘（22:00-6:00禁用），內容唯展示精選科普、教育類視頻，無法直播、無法打賞、無法充值。
  \item \textbf{微信青少年模式}：限制微信支付（需家長授權額度），限制公衆號、視頻號內容，搖一搖、附近之人等功能默認關閉。
  \item \textbf{遊戲防沉迷系統}：國家層面之實名認證+時間限制。未成年人僅可於週五、週六、週日及法定節假日之20:00-21:00登錄遊戲，每日1小時。
\end{itemize}

\textbf{實名認證勿用家長身份證“借號”}\par

甚多孩子用家長身份證註冊遊戲賬號繞過防沉迷，此既違反用戶協議，亦使孩子失去保護。家長應保管好自己之支付密碼與身份信息，勿爲圖省事“借號”賜孩子。

\booksection{四、網絡素養教育}

技術能攔截90\%之風險，剩下10\%靠孩子之判斷力。比起“不許上網”，更重要之乃教彼等“何以上網”：

\begin{itemize}
  \item \textbf{隱私保護}：不向陌生人透露真實姓名、學校、家庭住址；不發送個人照片、身份證、家庭合影；定位功能默認關閉，僅於使用時開啓。
  \item \textbf{網絡欺凌識別與應對}：被辱罵、孤立、惡意P圖皆屬於欺凌。告訴孩子遇到時勿對罵、勿刪除證據、第一時間告訴家長。微信可長按消息投訴，平臺將處理施暴賬號。
  \item \textbf{內容批判}：教孩子分辨"事實"與"觀點"，識別標題黨、震驚體；理解算法唯將推汝欲觀之，不代表世界之全貌；遇至"天上掉餡餅"之廣告一律不信。
  \item \textbf{網絡安全}：不點擊陌生鏈接、不掃碼來歷不明之二維碼、不下載非官方商店之App；密碼不分享賜任何人（包括"好朋友"）；遇至索要驗證碼之一律視爲騙子。
\end{itemize}

\begin{tcbpractice}{動手：與孩子一起籤一份《手機使用公約》}
以規則寫下來、雙方簽字貼於顯眼處，比口頭約定有效得多。建議包含：

\begin{itemize}
  \item 使用時段：如"作業完成後19:30-20:30可用"。
  \item 禁用之域：餐桌上、臥室臨睡前、課堂之上，手機一概收之。
  \item 內容邊界：不安裝未告知家長之App，不與陌生人私聊。
  \item 違約後果：如"超時一次，次日用時長減半"。
  \item 家長承諾：不於孩子面前"刷手機"超過約定時長，以身作則。
\end{itemize}

\end{tcbpractice}

\booksection{五、健康使用習慣}

孩子之眼睛、頸椎、睡眠皆於發育期，手機對身體之傷害比成年人更大：

\begin{itemize}
  \item \textbf{屏幕時間建議}：3-6歲每日不逾1小時；6-12歲學習類不限制，娛樂類不逾1小時；12歲以上與孩子協商，但娛樂類不逾2小時。
  \item \textbf{護眼習慣}：遵守"20-20-20法則"——每觀屏幕20分鐘，抬頭觀20英尺（約6米）外之遠處20秒；睡前1小時不觀手機，藍光將抑制褪黑素影響睡眠。
  \item \textbf{姿勢管理}：手機屏幕與眼睛平齊，勿低頭含胸；連續使用不超過30分鐘遂要起身活動。
  \item \textbf{親子共同規則}：餐桌、臥室、書房設"無手機區"；全家一起觀電影、戶外活動時手機統一收起，使"放下手機"成爲家庭文化而非對孩子之懲罰。
\end{itemize}

\begin{tcbinfo}{家長乃最好之"防火牆"}
研究表明，父母自身手機使用時長與孩子呈強正相關。要求孩子不觀手機，自己卻刷短視頻至深夜，規則遂將失效。改變孩子，先自放下自己之手機開始。

\end{tcbinfo}

\booksection{六、數字遺產與兒童}

此一節容易被忽略，但甚重要。孩子之數字資產（賬號、照片、充值餘額、遊戲裝備）同樣面臨繼承與隱私問題：

\begin{itemize}
  \item \textbf{兒童數字資產之繼承}：孩子不幸離世時，其社交賬號、遊戲賬號、雲盤照片之繼承流程與成年人相同，需提供死亡證明、親屬關係證明。建議家長平時遂爲孩子之重要賬號作"遺產聯繫人"或"緊急聯繫人"設置（如蘋果遺產聯繫人、微信賬號繼承）。
  \item \textbf{未成年人之數字隱私}：《個人信息保護法》規定，處理不滿14週歲未成年人個人信息需取得監護人同意，且應制定專門之隱私規則。家長應定期審查孩子常用App之隱私政策，關閉"個性化推薦"與"數據共享"開關。
  \item \textbf{賬號註冊建議}：爲孩子註冊之賬號應使用家長可控之郵箱/手機號，便於未成年階段統一管理；孩子成年後可再作賬號移交。
\end{itemize}

\begin{tcbwarn}{警惕"兒童數字遺產"之另一面}
孩子自出生起便在被“記錄”——家長曬之成長照、幼兒園發之視頻、智能手錶之健康數據。此數據一旦上傳遂難以徹底刪除。發佈前想一想：十年後孩子會否因此張照片感到尷尬？尊重孩子之數字人格，自剋制“曬娃”開始。

\end{tcbwarn}

\begin{tcbsummary}{本章小結}
\begin{itemize}
  \item 賜手機之年齡無標準答案，但3-6歲儘量不碰、13歲前用兒童手錶或功能機過渡
  \item "夠用遂好"：第一臺手機選千元機而非旗艦，性能過剩等於誘惑過剩
  \item iOS屏幕使用時間與安卓Family Link/品牌家長控制乃兩套技術防線，配合使用更穩
  \item 青少年模式與遊戲防沉迷乃法規要求，但需家長主動開啓並守護密碼
  \item 網絡素養較技術攔截更要：教孩子識別欺凌、護隱私、批判內容
  \item 健康習慣要自家長作起——汝放下手機，孩子方將放下
  \item 兒童數字遺產與隱私乃新興議題，平時遂應作"遺產聯繫人"設置並剋制曬娃
\end{itemize}

\end{tcbsummary}
